% !TeX root = ../navier-stokes-manuscript.tex
% ============================================================
% 3.3 The Critical Height at the Terminal Face
% ============================================================
\subsection{The Critical Height at the Terminal Face}\label{sub:TheCriticalHeightAtTheTerminalFace}

The whole-terminal rectangles contain adjacent spatial rows on \(I_*\).
Before using them at critical height, we prove that one such pair gives a strong value at \(T_*\).

\begin{lemma}[Adjacent Hilbert-scale trace]\label{lem:AdjacentHilbertScaleTrace}
Let \(s\in\R\), and let \(v\) satisfy
\[
  v\in L^2(I_*;H^{s+1}),
  \qquad
  v_t\in L^2(I_*;H^{s-1}).
\]
Then \(v\) has a unique representative in \(C([0,T_*];H^s)\).
For \(0\leq a\leq b\leq T_*\), this representative satisfies
\begin{equation}\label{eq:AdjacentHilbertScaleEstimate}
  \left|
  \norm{v(b)}_{H^s}^{2}
  -\norm{v(a)}_{H^s}^{2}
  \right|
  \leq
  2\norm{v_t}_{L^2(a,b;H^{s-1})}
  \norm{v}_{L^2(a,b;H^{s+1})}.
\end{equation}
\end{lemma}

\begin{proof}
Put \(A=(1-\Delta)^{1/2}\).
Time mollification and Fourier truncation first give, for smooth approximants,
\begin{equation}\label{eq:AdjacentHilbertScaleIdentity}
  \norm{v(b)}_{H^s}^{2}
  -\norm{v(a)}_{H^s}^{2}
  =
  2\operatorname{Re}\int_a^b
  \pair{A^{s-1}v_t(t)}{A^{s+1}v(t)}_{L^2}
  \,dt.
\end{equation}
Cauchy's inequality gives \eqref{eq:AdjacentHilbertScaleEstimate} for the approximants.
The two assumed \(L^2\) norms let the approximants converge in the right side of \eqref{eq:AdjacentHilbertScaleIdentity}.
The resulting identity makes the \(H^s\) norm continuous, while the same approximation gives weak \(H^s\) continuity.
Weak continuity together with norm continuity gives strong \(H^s\) continuity.
Two continuous representatives that agree almost everywhere agree everywhere, which proves uniqueness.
\end{proof}

The critical readout uses only the rectangle \(\cR_{0,2}[u_0]\).
That rectangle supplies \(u\in L^2(I_*;H^3)\) and \(u_t\in L^2(I_*;H^1)\), hence
\begin{equation}\label{eq:SectionThreeCriticalAdjacentPair}
  u\in L^2(I_*;H^{5/2}),
  \qquad
  u_t\in L^2(I_*;H^{1/2}).
\end{equation}
The fractional pair follows from the continuous inhomogeneous Sobolev embeddings, so no fractional-index rectangle is being assumed.
Applying the adjacent trace lemma with \(s=3/2\) gives
\begin{equation}\label{eq:SectionThreeCriticalVelocityTrace}
  u\in C([0,T_*];H^{3/2}).
\end{equation}

Let \(H(t)=E_{1/2}(t)\), with \(E_s\) defined in \eqref{eq:SectionThreeSpatialEnergy}.
The homogeneous identity needed here follows directly on the critical pair.
Apply time mollification together with Fourier multipliers that truncate to
\(\varepsilon\leq|\xi|\leq R\).
For the resulting smooth, frequency-localized functions, differentiation and Plancherel's theorem give
\[
  \frac{1}{2}\norm{\Lambda^{3/2}u(b)}_{L^2}^{2}
  -\frac{1}{2}\norm{\Lambda^{3/2}u(a)}_{L^2}^{2}
  =
  \operatorname{Re}\int_a^b
  \pair{\Lambda^{1/2}u_t(t)}{\Lambda^{5/2}u(t)}_{L^2}
  \,dt.
\]
The two memberships in \eqref{eq:SectionThreeCriticalAdjacentPair} give convergence of the integral by Cauchy's inequality.
The continuity in \eqref{eq:SectionThreeCriticalVelocityTrace} gives convergence of the endpoint terms as \(\varepsilon\downarrow0\), \(R\uparrow\infty\), and the time mollifier is removed.
Consequently, for \(0\leq a\leq b\leq T_*\),
\[
  E_{3/2}(b)-E_{3/2}(a)
  =
  \operatorname{Re}\int_a^b
  \pair{\Lambda^{1/2}u_t(t)}{\Lambda^{5/2}u(t)}_{L^2}
  \,dt.
\]
The pair in \eqref{eq:SectionThreeCriticalAdjacentPair} places the integrand in \(L^1(I_*)\).
Equation \eqref{eq:SectionThreeAllSpatialEnergies} supplies \(E_{3/2}\in L^1(I_*)\).
Thus
\begin{equation}\label{eq:SectionThreeCriticalEnergyRegularity}
  E_{3/2}\in W^{1,1}(I_*)\cap L^1(I_*)
  \subset L^\infty(I_*).
\end{equation}

The stronger integer rows \(u,u_t\in L^2(I_*;H^1)\) give
\begin{equation}\label{eq:SectionThreeCriticalHeightDerivative}
  H'(t)
  =\pair{\Lambda^{1/2}u(t)}{\Lambda^{1/2}u_t(t)}_{L^2}
  \in L^1(I_*).
\end{equation}
The pressure-explicit critical production current is
\[
  \mathcal P_{1/2}
  :=
  -\operatorname{Re}
  \pair{\Lambda^{1/2}u}
  {\Lambda^{1/2}\bigl((u\cdot\nabla)u+\nabla p\bigr)}_{L^2}.
\]
Pairing the original equation with \(\Lambda u\) yields the exact balance
\[
  H'=-2\nu E_{3/2}+\mathcal P_{1/2}.
\]
Equations \eqref{eq:SectionThreeCriticalEnergyRegularity} and \eqref{eq:SectionThreeCriticalHeightDerivative} now give
\[
  H\in W^{1,1}(I_*),
  \qquad
  E_{3/2}\in L^\infty(I_*)\cap L^1(I_*),
  \qquad
  \mathcal P_{1/2}\in L^1(I_*).
\]
Thus the critical height and \(E_{3/2}\) have continuous representatives up to \(T_*\), while the pressure-explicit production current is integrable in time.
No endpoint value is assigned to that current.
The rectangle \(\cR_{0,2}[u_0]\) supplies this critical trace; the higher rectangles will now give one terminal velocity in every finite Sobolev space.
